\section{Household Optimality Conditions}
\label{app:household}

In this Appendix, I derive the conditions that Section~\ref{sec:m-analysis}
uses: the choices within a housing option, the Euler equation for the bond
and the borrowing multiplier $\zeta_t$, the effect of credit on the value of
a household and on fertility, and the two results in the footnotes of that
Section. Throughout, $n$ is the number of children at home after the birth
outcome of period $t$, and the budget constraint \eqref{eq:m-budget} of the
chosen option defines $c_t$ as a function of $b_{t+1}$ and, for a renter, of
$h^R_t$.

\subsubsection*{Choices within an option.}
Within option $j$ of \eqref{eq:m-vh}, renting or owning $h^O_t\in\mathcal H$,
the household chooses $b_{t+1}$ and, if it rents, $h^R_t$, subject to the
borrowing limit \eqref{eq:m-ltv}, with $h^O_t=0$ for a renter, and to the
cap $h^R_t\le\bar h^R$. Let $\zeta_t\ge0$ and $\nu_t\ge0$ be the multipliers
on the two constraints. By \eqref{eq:m-budget},
$\partial c_t/\partial b_{t+1}=-1$ in every option and
$\partial c_t/\partial h^R_t=-r_t$ for a renter. In \eqref{eq:m-utility}, the
marginal utilities of consumption and of rented space are the same power of
the composite in brackets, times $\alpha_0/c_t$ and
$(1-\alpha_0)/(h^R_t-\bar h(n))$, so the condition
$\partial u/\partial h^R_t=r_t\,\partial u/\partial c_t+\nu_t$ gives, below
the cap,
\begin{equation}
 r_t\bigl(h^R_t-\bar h(n)\bigr)=\frac{1-\alpha_0}{\alpha_0}\,c_t .
 \label{eq:app-intra}
\end{equation}
A renter spends the share $\alpha_0$ of $c_t+r_t(h^R_t-\bar h(n))$ on
consumption and $1-\alpha_0$ on space above the requirement. With
$b_t=b_{t+1}=0$ and no home to sell, this spending is
$y_t+T_t-r_t\bar h(n)$, and substituting the two shares into
\eqref{eq:m-utility} gives \eqref{eq:m-indirect}. At the cap,
$h^R_t=\bar h^R$ and
$\nu_t=\partial u/\partial h^R_t-r_t\,\partial u/\partial c_t\ge0$. An
owner's space is fixed within the option, so its only choice is the bond.

\subsubsection*{The bond.}
The first-order condition for $b_{t+1}$ is the Euler equation
\begin{equation}
 \frac{\partial u}{\partial c_t}
 =\beta\,\E_t\Bigl[s_a\frac{\partial V_{t+1}}{\partial b_{t+1}}+(1-s_a)\,B'(x_{t+1})\Bigr]+\zeta_t,
 \qquad
 \zeta_t\bigl(b_{t+1}+\phi P_th^O_t\bigr)=0,
 \label{eq:app-euler}
\end{equation}
where $h^O_t=0$ for a renter and
$B'(x_{t+1})=\theta_0(\theta_1+x_{t+1})^{-\sigma}$. The multiplier $\zeta_t$
is the gain, in period-$t$ utility, from borrowing one more unit in period
$t$ and repaying $R_b$ in period $t+1$: the marginal utility of consumption
today net of the discounted expected marginal value of the bond tomorrow. It
is positive only at the limit. The estate $x_{t+1}=b_{t+1}+(1-\mu)P_th^O_t$
contains the bond position itself, so $\partial x_{t+1}/\partial b_{t+1}=1$
and the bequest term carries no $R_b$; the survival term does, through the
budget of period $t+1$.

For Gumbel shocks with scale $\kappa$, the expected maximum over options is
the logsum in \eqref{eq:m-fertility} and \eqref{eq:m-tenure}, whose
derivative with respect to the value of one option is the probability of
that option:
\begin{equation}
 \frac{\partial}{\partial V^k}\,\kappa\log\sum_{j}\exp\bigl(V^j/\kappa\bigr)
 =\frac{\exp\bigl(V^k/\kappa\bigr)}{\sum_{j}\exp\bigl(V^j/\kappa\bigr)} .
 \label{eq:app-logsum}
\end{equation}
The household observes the shocks before it chooses, so a marginal change in
a state moves the value of the chosen option only; at the shocks where the
choice switches, the two options have the same value. The derivative of the
expected maximum with respect to a state is therefore the
probability-weighted average of the derivatives of the option values. When
$\sigma\ge1$, the value of an option whose feasible set shrinks to a point
tends to $-\infty$, since $u\to-\infty$ as $c_t\to0$, so its probability
tends to zero and the logsum is continuous where an option enters or leaves
the feasible set. Within an option $j$ of period $t+1$, $b_{t+1}$ enters
only the budget, with $\partial c_{t+1}/\partial b_{t+1}=R_b$, so
$\partial V^j_{t+1}/\partial b_{t+1}=R_b\,\partial u/\partial c_{t+1}$ at the
option's consumption, whether or not a constraint binds in that period.
Writing $\frac{\partial u}{\partial c_{t+1}}(j,n')$ for this marginal
utility with $n'$ children at home, and applying \eqref{eq:app-logsum} to
\eqref{eq:m-tenure} and to \eqref{eq:m-fertility},
\begin{multline}
 \frac{\partial V_{t+1}}{\partial b_{t+1}}
 =R_b\Bigl[\bigl(1-\pi_{a+1}\pi^F_{t+1,a+1}\bigr)\sum_{j}\pi^H_{t+1}(j\mid n_{t+1})\,\frac{\partial u}{\partial c_{t+1}}(j,n_{t+1})\\
 +\pi_{a+1}\pi^F_{t+1,a+1}\sum_{j}\pi^H_{t+1}(j\mid n_{t+1}+1)\,\frac{\partial u}{\partial c_{t+1}}(j,n_{t+1}+1)\Bigr],
 \label{eq:app-envelope}
\end{multline}
where $n_{t+1}$ is the number of children at home at the start of period
$t+1$ and $\pi_{a+1}\pi^F_{t+1,a+1}$ is the probability of a birth in that
period, zero outside the fertile ages. The derivative is $R_b$ times the
marginal utility of consumption in period $t+1$, averaged over the birth
outcome and the housing choice of that period; the cost $\xi$ drops out. The
expectation $\E_t$ in \eqref{eq:app-euler} then averages over $z_{t+1}$ and
the departure of children.

\subsubsection*{A birth and the multiplier.}
At the limit, $\zeta_t$ is the difference between the two sides of
\eqref{eq:app-euler} at $b_{t+1}=-\phi P_th^O_t$, and a birth changes both.
In period $t$, it adds $r_th_P$ to the spending a renter cannot avoid,
lowers an owner's space above the requirement by $\chi h_P$, and raises
$e(n)$; for $\sigma>1$, each of these raises $\partial u/\partial c_t$ at the
given bond position. In period $t+1$ the same holds with probability
$1-\lambda$, when the child is still at home, and the bequest term does not
depend on $n$. Within an option,
\begin{equation}
 \zeta_t(h^O_t,n+1)-\zeta_t(h^O_t,n)
 =\Bigl[\frac{\partial u}{\partial c_t}\Bigr]_{n}^{n+1}
 -\beta s_a\,\E_t\Bigl[\frac{\partial V_{t+1}}{\partial b_{t+1}}\Bigr]_{n}^{n+1},
 \label{eq:app-birth}
\end{equation}
where $[\,\cdot\,]_n^{n+1}$ is the difference between the values with $n+1$
and with $n$ children at home, at $b_{t+1}=-\phi P_th^O_t$ in both. The
first term is positive, and larger when spending in period $t$ is small,
since the marginal utility of consumption is convex in spending. The second
is scaled by $\beta s_a$ and by $1-\lambda$, averaged with the choice
probabilities of period $t+1$, and evaluated at that period's spending,
which is larger when earnings rise with age. The multiplier therefore rises
for a young household with low spending today and rising earnings, and can
fall for one that expects to spend less next period; the sign is ambiguous
in general.

\subsubsection*{Credit.}
Consider a change in $\phi$ in period $t$ alone. It enters the problem of an
owner only through the borrowing limit, whose derivative with respect to
$\phi$ is $P_th^O_t$, so by the envelope theorem
$\partial V^j_t/\partial\phi=\zeta_t(h^O,n)P_th^O$ for owning
$h^O\in\mathcal H$ and $\partial V^j_t/\partial\phi=0$ for renting, where
$\zeta_t(h^O,n)$ is the multiplier in the option of owning $h^O$ with $n$
children at home. Applying \eqref{eq:app-logsum} to \eqref{eq:m-tenure}
gives the second line of \eqref{eq:m-credit}. For the probability of trying,
\eqref{eq:m-fertility} and
$V^F_{t,A}-V^F_{t,0}=\pi_a\bigl[V^H_t(n_t+1)-V^H_t(n_t)-\xi\,\1\{\text{first child}\}\bigr]$
give
\begin{equation}
 \frac{\partial\pi^F_{t,a}}{\partial\phi}
 =\frac{\pi^F_{t,a}\bigl(1-\pi^F_{t,a}\bigr)}{\kappa}\,
 \frac{\partial\bigl(V^F_{t,A}-V^F_{t,0}\bigr)}{\partial\phi}
 =\frac{\pi^F_{t,a}\bigl(1-\pi^F_{t,a}\bigr)}{\kappa}\,\pi_a
 \left[\frac{\partial V^H_t(n_t+1)}{\partial\phi}-\frac{\partial V^H_t(n_t)}{\partial\phi}\right],
 \label{eq:app-credit-pi}
\end{equation}
where $\kappa=\kappa_1$ before the first child and $\kappa_C$ after; the
cost $\xi$ does not depend on $\phi$ and drops out. This is the first line
of \eqref{eq:m-credit}. If the change in $\phi$ is permanent, it also enters
$V_{t+1}$, and the envelope theorem adds
$\beta s_a\E_t\,\partial V_{t+1}/\partial\phi$ to the derivative of each
option, which by \eqref{eq:app-logsum} is the second line of
\eqref{eq:m-credit} in period $t+1$, averaged over the choices of that
period, and so on: the derivative becomes the sum of these terms over all
future periods, discounted by $\beta$ and weighted by survival and by the
probabilities of the choices that lead to each option.

The bracket in \eqref{eq:m-credit} splits into two parts,
\begin{multline}
 \frac{\partial V^H_t(n+1)}{\partial\phi}-\frac{\partial V^H_t(n)}{\partial\phi}
 =\sum_{h^O\in\mathcal H}\bigl[\pi^H_t(h^O\mid n+1)-\pi^H_t(h^O\mid n)\bigr]\,\zeta_t(h^O,n)\,P_th^O\\
 +\sum_{h^O\in\mathcal H}\pi^H_t(h^O\mid n+1)\bigl[\zeta_t(h^O,n+1)-\zeta_t(h^O,n)\bigr]\,P_th^O ,
 \label{eq:app-credit-decomp}
\end{multline}
where the first sum holds the multipliers at their values without the child
and the second holds the probabilities at their values with the child. The
first sum is positive when the child moves probability from renting, and
from homes where the limit is slack, toward homes where
$\zeta_t(h^O,n)P_th^O$ is large. The second is positive when the child
raises the multiplier in the homes the household is likely to own with it,
with the sign of \eqref{eq:app-birth}. Both are small for a household that
rents with high probability with and without the child.

\subsubsection*{The owner's cost of space.}
By \eqref{eq:m-budget}, an owner that holds a unit of housing from period $t$
to period $t+1$ pays $P_t$ and $\tau^pP_t$ in period $t$. In period $t+1$ it
pays maintenance $\delta_HP_{t+1}$ if it stays, or receives
$(1-\delta_H-\mu)P_{t+1}$ if it sells, so apart from the selling cost the unit
is worth $(1-\delta_H)P_{t+1}$ in period $t+1$, which the bond values at
$(1-\delta_H)P_{t+1}/R_b$ in period $t$. A renter pays $r_t$ in period $t$. By
\eqref{eq:m-rent},
\begin{equation}
 r_t-\Bigl[(1+\tau^p)P_t-\frac{(1-\delta_H)P_{t+1}}{R_b}\Bigr]=0,
 \label{eq:app-usercost}
\end{equation}
where the term in brackets is the owner's cost in period-$t$ units. Apart from
the selling cost, owning a unit of space costs the same as renting it, along a
transition as well as in a stationary equilibrium.

\subsubsection*{The cost of a child's space and earnings.}
Let $\tilde y_t=y_t+T_t$, $k=r_th_P$ and
$\Theta_t=\alpha_0^{\alpha_0}(1-\alpha_0)^{1-\alpha_0}r_t^{\alpha_0-1}$. The
first term of \eqref{eq:m-indirect} is $(\Theta_t\tilde y_t)^{1-\sigma}/(1-\sigma)$
without children and
$\bigl(\Theta_t(\tilde y_t-k)/e(1)\bigr)^{1-\sigma}/(1-\sigma)$ with one
child, so the utility cost of a first child and its second derivative are
\begin{equation}
\begin{aligned}
 \Delta_t(\tilde y_t)&=\frac{\Theta_t^{1-\sigma}}{1-\sigma}\Bigl[\tilde y_t^{\,1-\sigma}-e(1)^{\sigma-1}(\tilde y_t-k)^{1-\sigma}\Bigr],\\
 \Delta_t''(\tilde y_t)&=\sigma\,\Theta_t^{1-\sigma}\Bigl[e(1)^{\sigma-1}(\tilde y_t-k)^{-\sigma-1}-\tilde y_t^{\,-\sigma-1}\Bigr],
\end{aligned}
\label{eq:app-delta}
\end{equation}
where $\tilde y_t>k$. For $\sigma\ge1$, $e(1)^{\sigma-1}\ge1$ and
$(\tilde y_t-k)^{-\sigma}>\tilde y_t^{\,-\sigma}$, so
$\Delta_t'(\tilde y_t)=\Theta_t^{1-\sigma}\bigl[\tilde y_t^{\,-\sigma}-e(1)^{\sigma-1}(\tilde y_t-k)^{-\sigma}\bigr]<0$
and $\Delta_t''>0$: the cost is decreasing and convex in earnings. The
second derivative rises with $k$, since
$\partial\Delta_t''/\partial k=\sigma(\sigma+1)\Theta_t^{1-\sigma}e(1)^{\sigma-1}(\tilde y_t-k)^{-\sigma-2}>0$,
and, for $\sigma>1$, with $r_t$, since $\Theta_t^{1-\sigma}$ is proportional
to $r_t^{(1-\alpha_0)(\sigma-1)}$. For a renter at the cap with and without
the child, $c_t=\tilde y_t-r_t\bar h^R$ in both cases and the child lowers
space above the requirement from $\bar h^R$ to $\bar h^R-h_P$, so
\[
 \Delta_t(\tilde y_t)=\frac{(\tilde y_t-r_t\bar h^R)^{\alpha_0(1-\sigma)}}{1-\sigma}
 \Bigl[(\bar h^R)^{(1-\alpha_0)(1-\sigma)}-e(1)^{\sigma-1}\bigl(\bar h^R-h_P\bigr)^{(1-\alpha_0)(1-\sigma)}\Bigr],
\]
which for $\sigma>1$ is positive, decreasing and convex in $\tilde y_t$ and
rises with $r_t$ and $h_P$.
